\documentclass[letterpaper, 11pt]{article}

\usepackage{
	hyperref,
	enumitem,
	longtable,
	fancyhdr,
	amsmath,
	array,
	libertine,
	newtxmath
}
\usepackage[dvipsnames]{xcolor}

\usepackage[
	left=0.7in,
	right=0.8in,
	top=0.8in,
	bottom=0.8in,
	headsep=0in,
	footskip=0.2in
]{geometry}

\hypersetup{
	colorlinks=true,
	linkcolor=blue,
	urlcolor=blue,
	pdfborderstyle={/S/U/W 1},
	linkbordercolor=blue,
}

\renewcommand{\baselinestretch}{1.2}
\pagenumbering{gobble}

\newcommand{\updateinfo}[1][\today]{\par
\vfill
\hfill{\scriptsize\color{gray}Updated #1}}
\newcommand{\headingfont}{\Large
\color{Black}}
\newcolumntype{R}{>{\raggedleft}p{1in}}

\newcommand{\SectionHeading}[2][0pt]{ & {\headingfont\textbf{#2}}\\[#1] }

\newenvironment{SectionTableBase}[1]{%
\renewcommand{\arraystretch}{#1}%
\setlength{\tabcolsep}{10pt}%
\begin{longtable}{Rp{5.2in}} }{%
\end{longtable}
\vspace{-.3cm}
}

\newenvironment{SectionTable}{\begin{SectionTableBase}{1.7}}{\end{SectionTableBase}}
\newenvironment{SectionTableSingleSpace}{\begin{SectionTableBase}{1.2}}{\end{SectionTableBase}}
\newenvironment{SectionTableHalfSpace}{\begin{SectionTableBase}{0.9}}{\end{SectionTableBase}}
% Setup fancyhdr for the first page
\fancypagestyle{firstpage}{ \fancyhf{} \fancyfoot[R]{\updateinfo} \renewcommand{\headrulewidth}{0pt} \renewcommand{\footrulewidth}{0pt} }

% Apply the firstpage style to the first page
\AtBeginDocument{ \thispagestyle{firstpage} }

\begin{document}
\begin{SectionTable}
    & {\Huge Tiwa Eisape} \\[0pt] & Machine learning, Natural Language Processing, Cognitive Science
    \newline
    \href{mailto:te1187@princeton.edu}{te1187@princeton.edu} $\;|\;$ \href{https://eisa.pe}{https://eisa.pe}
\end{SectionTable}

\begin{SectionTableSingleSpace}
    \SectionHeading{Academic Positions} %
    2025 -- & Postdoctoral Researcher, Princeton University
    \newline
    Advisors: Tom Griffiths, Adele Goldberg \\%
    2022 -- 2025 & Student Board Member, Cognitive Science Society
\end{SectionTableSingleSpace}

\begin{SectionTableSingleSpace}
    \SectionHeading{Education} %
    2025 & Ph.D, Cognitive Science, MIT %
    \newline
    Advisors: Roger Levy, Yoon Kim \\%
    2019 & B.S, Computer Science with Minor in Philosophy, Boston College%
    \newline
    Advisors: Joshua Hartshorne, Stefano Anzellotti%
\end{SectionTableSingleSpace}



\begin{SectionTableSingleSpace}
    \SectionHeading{Research Positions}%
    2024 & FAIR Labs (Meta AI), Research Scientist Intern%
    \newline
    Supervisors: Adina Williams, Karen Ullrich \\%
    2023 & Google AI, Student Researcher%
    \newline
    Supervisors: Tal Linzen, Fei Sha \\%
    2019 & Vimeo Machine Learning, Research Scientist Intern
\end{SectionTableSingleSpace}

\begin{SectionTableSingleSpace}
    \SectionHeading{Honors \& Awards}%
    2025 & Presidential Postdoctoral Fellowship (Princeton) \\%
    2021 & Graduate Research Fellowship (NSF) \\%
    2019 & GEM Full Fellowship (InterActiveCorp/GEM Consortium) \\%
    2016 & Ronald E. McNair Scholarship (TRIO) \\%
    2019 & Dean of Sciences Fellowship (MIT) \\%
    2019 & Presidential Fellowship (MIT)%
    \newline
    \textit{awarded \& declined} \\%
    2019 & Matthew R. Copithorne Scholarship (Boston College) \\%
    2021 & Walle Nauta Teaching Award (MIT) \\%
    2019 & Order of the Cross and Crown (Boston College) \\%
\end{SectionTableSingleSpace}

% \newpage

\begin{SectionTableHalfSpace}
    \SectionHeading{Advising} %
    2022 -- 2023 & Jesus Crespo, MIT,M.Eng '23, Computation and Cognition,
    MIT\\%
    &\textsl{Head-Seeking RNNGs: On the Relationship between Top-down Prediction
        and Evidence Integration in Parsing.} \\%
    2023 & Zhe Su, Machine Learning (CMU) \\%
    2022 -- 2023 & Vineet Gangireddy, Applied Mathematics (Harvard): \textsl{A
        Computational Approach to Recontextualization in Human Reading Behavior.}
    \\%
    2022 & Nicole Wong, Computer Science (MIT) \\ %
    2022 & Heidi Lei, Mathematics (MIT) \\%
    2021 & Vinh Le, Computation and Cognition (MIT) \\%
    2020 & Robert Chen, Computer Science (MIT) \\%
    2020 & Carl Joshua Quines, Mathematics (MIT)
\end{SectionTableHalfSpace}

\begin{SectionTableSingleSpace}
    \SectionHeading{Professional Service}

    2024 & Founding Organizer of the Graduate Student Steering Committee of the
    Cognitive Science Society \\

    2022 -- & Diversity \& Inclusion Committee, Cognitive Science Society \\

    2022 -- & Membership Committee, Cognitive Science Society \\

    2022 & Organizer, MIT CompLang \\

    2019 & Organizer, Vimeo Machine Learning seminar \\

    2021 & Technical Session Chair, Conference of the Cognitive Science Society
    \\

    2020 -- & Diversity Ambassador, Office of Graduate Education (MIT) \\

    2022 & MIT Summer Research Program Mentor, Office of Graduate Education
    (MIT)\\

    2020 & Reviewer, DataKind/Data.org Inclusive Growth \& Recovery Challenge\\

    2017 & Tutor, Resilient Coders \\

    2018 & Preceptor, Options Through Education (Boston College) \\

    2017 & Council Member, PULSE Program for Service Learning (Boston College)
    \\

    & \textbf{Reviewing (journals)}: OpenMind ('21, '23) \\

    & \textbf{Reviewing (conferences)}: CogSci ('20, '21, '23, '24)
\end{SectionTableSingleSpace}

\begin{SectionTableHalfSpace}
    \SectionHeading{Talks}

    2024 & Brain \& AI (PI: Jean Remi-King), Meta AI \\

    2024 & Reasoning, Attention, \& Memory group (PI: Jason Weston), Meta AI
    \\

    2023 & Stanford NLP Seminar, Stanford \\

    2021 & How Can Findings about the Brain Improve AI Systems workshop,
    ICLR\\

    2021 & Wehbe Lab, Carnegie Mellon University \\

    2021 & EdinburghNLP, The University of Edinburgh \\

    2021 & Psycholinguistics Coffee, The University of Edinburgh \\

    2019 & Machine Learning Reading Group, Vimeo \\
\end{SectionTableHalfSpace}

\begin{SectionTableHalfSpace}
    \SectionHeading{Teaching} & \textbf{Graduate courses} \\

    (Spring) 2021 & Teaching Assistant, Cognitive Science (MIT) \\

    (Fall) 2020 & Teaching Assistant, Computational Cognitive Science (MIT) \\

    (Fall) 2018 & Teaching Assistant, Computational Models of Cognition (Boston
    College) \\

    & \textbf{Undergraduate courses} \\

    (Fall) 2018 & Teaching Assistant, Data Mining (Boston College) \\

    (Spring) 2019 & Teaching Assistant, Algorithms (Boston College) \\

    (Spring) 2018 & Teaching Assistant, Computer Science I (Boston College) \\

    (Spring) 2018 & Teaching Assistant, Personal and Social Responsibility (Boston
    College) \\

    (Fall) 2018 & Teaching Assistant, Personal and Social Responsibility (Boston
    College) \\

    (Fall) 2017 & Teaching Assistant, Personal and Social Responsibility (Boston
    College)
\end{SectionTableHalfSpace}

\begin{SectionTable}
    \SectionHeading[-1em]{Publications}%
    & Undergraduate and Master's Advisees are \underline{underlined} \\%
    &\textbf{\underline{Thesis}}\\& \textbf{Eisape, T.} (2025). Language Comprehension, Production, and Reasoning in Humans and Neural Language Models. Massachusetts Institute of Technology.\\%
    &\textbf{\underline{Papers}}\\
    & \textbf{Eisape, T.}, Tessler, M., Dasgupta, I., Sha, F., Steenkiste, S.,
    \& Linzen, T. (2024). A Systematic Comparison of Syllogistic Reasoning
    in Humans and Language Models. NAACL.\\%
    & Zhou, X., \underline{Su, Z.}, \textbf{Eisape, T.}, Kim, H., \& Sap, M.
    (2024). Is this the real life? Is this just fantasy? The Misleading Success
    of Simulating Social Interactions With LLMs. EMNLP.\\%
    & \textbf{Eisape, T.}, \underline{Gangireddy, V.}, Levy, R. \& Kim, Y. (2022).
    Probing for Incremental Parse States in Autoregressive Language Models. EMNLP.\\%
    & Liu, Z., \textbf{Eisape, T.}, Prud'hommeaux, E., \& Hartshorne, J. (2022).
    Data-driven Crosslinguistic Syntactic Transfer in Second Language Learning.
    CogSci.\\%
    &Tucker, M., \textbf{Eisape, T.}, Qian, P., Levy, R., \& Shah, J. (2022).
    When Does Syntax Mediate Neural Language Model Performance? Evidence from
    Dropout Probes. NAACL.\\%
    &\underline{Chen, R.}, Levy, R., \& \textbf{Eisape, T.}(2021). On
    Factors Influencing Typing Time: Insights from a Viral Online Typing Game.
    CogSci/HSP.\\%
    &\textbf{Eisape, T.}, Zaslavsky, N., \& Levy, R. (2020). Cloze
    Distillation: Improving Neural Language Models with Human Next-Word Predictions.
    CoNLL.\\%
    &\textbf{Eisape, T.}, Levy, R., Tenenbaum, J., \& Zaslavsky, N. (2020). Toward
    human-like object naming in artificial neural systems. Bridging AI and
    Cognitive Science workshop, ICLR.\\%
    &\textbf{\underline{Abstracts \& Preprints}}\\&\textbf{Eisape, T.}, Merrill, W.,
    Diet z, S., \& Hartshorne, J. (2020). Grammatical Accents: Using Machine
    Learning to Quantify Language Transfer. CogSci/HSP.\\%
    & \textbf{Eisape, T.}, Grasso, S., Kurt, L., Lim, C., \& Hartshorne, J.
    (2022). Replication of Conway and Christiansen (2005) \, Modality-Constrained
    Statistical Learning of Tactile, Visual, and Auditory Sequences.\\%
    & Choong, H., \textbf{Eisape, T.}, \ldots \& Hartshorne, J. (2018). Fourth
    Replication of Saffran, Newport, \& Aslin (1996) Word segmentation: The role
    of distributional cues. Preprint.
\end{SectionTable}

\begin{SectionTableHalfSpace}
    \SectionHeading{Miscellaneous}

    2016 & Captain, Boston College Varsity Men's Rowing
\end{SectionTableHalfSpace}
\end{document}
% Eisape, T. (2025). Language Comprehension, Production, and Reasoning in Humans and Neural Language Models. Massachusetts Institute of Technology.
